% =====================================================================
%  CHƯƠNG 5: KẾT LUẬN VÀ HƯỚNG PHÁT TRIỂN
% =====================================================================
\chapter{Kết luận và hướng phát triển}
\label{chap:conclusion}

\section{Kết luận}

Đồ án đã xây dựng thành công một hệ thống gợi ý sản phẩm \textit{Software} trên
Amazon theo hướng tiếp cận hiện đại, kết hợp mô hình ngôn ngữ và kiến trúc hai tháp.
Các kết quả chính đạt được bao gồm:

\begin{itemize}[leftmargin=2em, itemsep=2pt]
    \item \textbf{Luồng xử lý dữ liệu lớn ổn định:} hoàn thiện pipeline end-to-end
    trên Apache Spark, từ nạp dữ liệu thô, làm sạch, ghép nối tới việc dựng văn bản
    ngữ nghĩa thống nhất và lưu trữ tối ưu bằng định dạng cột Parquet.
    \item \textbf{Thuật toán trọng số thời gian hiệu quả:} thiết kế cơ chế tổng hợp
    vector lịch sử người dùng kết hợp độ gần đây và mức độ ưa thích, tính toán hoàn
    toàn trên bộ nhớ RAM, nhẹ và nhanh.
    \item \textbf{Mô hình chứng minh năng lực truy hồi:} với hàm mất mát tương phản
    InfoNCE và bộ mã hóa ngôn ngữ làm tháp sản phẩm, mô hình đạt kết quả vượt trội
    hơn $24$ lần ở Top-$10$ và hơn $14$ lần ở Top-$50$ so với chọn ngẫu nhiên trên
    kho hơn $17.000$ sản phẩm, được kiểm chứng bằng giao thức đánh giá full-catalog
    khắt khe.
\end{itemize}

Qua đó, đồ án không chỉ giải quyết được hai hạn chế của hệ thống gợi ý truyền thống
(bỏ qua văn bản ngữ nghĩa và khởi động lạnh) mà còn thể hiện tư duy giải quyết bài
toán ở quy mô dữ liệu lớn -- đúng tinh thần của môn học.

\section{Hạn chế}

Bên cạnh các kết quả đạt được, hệ thống vẫn còn một số hạn chế:

\begin{itemize}[leftmargin=2em, itemsep=2pt]
    \item \textbf{Tháp sản phẩm bị đóng băng:} do giữ nguyên trọng số của mô hình
    ngôn ngữ để tiết kiệm bộ nhớ, tháp sản phẩm chưa được \textit{tinh chỉnh sâu}
    theo tư duy mua sắm thực tế của người dùng trong lĩnh vực công nghệ. Không gian
    biểu diễn sản phẩm vì thế chưa tối ưu, tạo ra một "trần" hiệu năng.
    \item \textbf{Chỉ số tuyệt đối còn thấp:} đặc thù của bài toán truy hồi trên kho
    rất lớn và dữ liệu thưa khiến Recall/NDCG tuyệt đối còn khiêm tốn, dù đã vượt xa
    baseline.
    \item \textbf{Chưa khai thác đặc trưng ngữ cảnh:} các yếu tố như mùa vụ, phiên
    truy cập hay quan hệ sản phẩm đi kèm chưa được tích hợp.
\end{itemize}

\section{Hướng phát triển}

Từ các hạn chế trên, đồ án đề xuất những hướng phát triển tiếp theo:

\begin{itemize}[leftmargin=2em, itemsep=2pt]
    \item \textbf{Tinh chỉnh tháp sản phẩm bằng kỹ thuật hiệu quả tham số:} khi có
    phần cứng mạnh hơn, áp dụng các kỹ thuật tinh chỉnh cục bộ như \textbf{LoRA}
    (Low-Rank Adaptation) để "mở băng" một phần tháp sản phẩm, giúp không gian biểu
    diễn thích ứng với hành vi người dùng mà vẫn tiết kiệm tài nguyên.
    \item \textbf{Tích hợp tìm kiếm vector siêu tốc:} sử dụng thư viện
    \textbf{FAISS} của Facebook để lập chỉ mục toàn bộ vector sản phẩm, cho phép hệ
    thống trả về gợi ý theo thời gian thực trong vài phần nghìn giây ngay cả với kho
    hàng triệu sản phẩm.
    \item \textbf{Mở rộng mô hình hành vi:} bổ sung mô hình hóa chuỗi (sequential)
    bằng các kiến trúc dựa trên Transformer (SASRec, BERT4Rec) cho tháp người dùng,
    cùng các tín hiệu ngữ cảnh thời gian thực.
    \item \textbf{Hệ thống lai và đa mục tiêu:} kết hợp truy hồi hai tháp ở tầng sinh
    ứng viên với một tầng xếp hạng tinh, đồng thời cân nhắc các tiêu chí đa dạng và
    độ mới của gợi ý.
\end{itemize}

\subsection{Phân tích hai hướng nâng cấp trọng tâm}

\textbf{Tinh chỉnh bằng LoRA.} Thay vì cập nhật toàn bộ $\sim 22{,}7$ triệu tham số
của tháp sản phẩm (tốn kém và dễ "quên" tri thức ngôn ngữ), LoRA \textit{đóng băng}
trọng số gốc và chèn thêm hai ma trận hạng thấp $A, B$ vào mỗi tầng:
$W' = W + \Delta W = W + BA$, với $A \in \mathbb{R}^{r \times d}$,
$B \in \mathbb{R}^{d \times r}$ và hạng $r \ll d$. Chỉ $A, B$ được huấn luyện -- số
tham số cập nhật giảm hàng trăm lần, cho phép "mở băng" tháp sản phẩm ngay cả trên
phần cứng hạn chế, giúp không gian biểu diễn thích ứng với hành vi mua sắm thực tế
mà vẫn giữ tri thức nền.

\textbf{Truy hồi thời gian thực với FAISS.} Ở giai đoạn suy luận, thay vì nhân ma
trận vét cạn, ta lập chỉ mục toàn bộ vector sản phẩm bằng FAISS. Với kho vừa phải có
thể dùng chỉ mục phẳng (\texttt{IndexFlatIP}) cho kết quả chính xác tuyệt đối; với
kho hàng triệu sản phẩm, các chỉ mục xấp xỉ như \texttt{IVF} hoặc \texttt{HNSW} cho
phép truy vấn Top-$N$ trong vài phần nghìn giây với độ chính xác gần như không suy
giảm -- đáp ứng yêu cầu thời gian thực của hệ thống sản xuất.

Tóm lại, đồ án đã đặt nền móng vững chắc cho một hệ thống gợi ý dựa trên ngữ nghĩa và
có khả năng mở rộng, sẵn sàng cho các cải tiến hướng tới môi trường sản xuất thực tế.
